\section{Statistical tests of geometric chaos}
\label{sec:statistics}

No single statistic is sufficient to identify chaos in a degenerate fiber.
We therefore use a hierarchy of tests.  Exact energy degeneracy supplies the
null result at the spectral level.  Curvature eigenvalues test local and
mesoscopic correlations of one geometric component.  Response trace words
then test whether the joint distribution of several physical tangent channels
is fixed by its complete two-point structure.  The hierarchy prevents a local
level-repulsion result from being promoted into a statement about Gaussian
closure or a thermodynamic limit.

\subsection{Spectral silence and curvature spectra}

For a protected multiplet, shift the common energy $E_0$ separately in each
sample and define
\[
 Z_E(t)=\Tr_P e^{-it(H-E_0)}.
\]
Exact degeneracy gives $Z_E(t)=D$.  With the normalization used here,
\begin{align}
 K_{E,\mathrm{raw}}(t)&=D^{-1}\mathbb E|Z_E(t)|^2=D,\nonumber\\
 K_{E,c}(t)&=D^{-1}\left(\mathbb E|Z_E(t)|^2
 -|\mathbb EZ_E(t)|^2\right)=0.
\end{align}
The equalities are algebraic and contain no dynamical inference.  In
particular, an energy ramp can be manufactured only by splitting the protected
multiplet or by mixing it with states outside the isolated fiber.

The curvature retains a nontrivial spectrum.  For a registered tangent pair
$(a,b)$, we remove its scalar part and fix its scale by
\[
 \widetilde F_{ab}=F_{ab}-D^{-1}(\Tr F_{ab})P,
 \qquad
 \sigma_{ab}^2=D^{-1}\Tr(\widetilde F_{ab}^{2}).
\]
Samples with $\sigma_{ab}=0$ are recorded as scalar controls rather than
silently normalized.  Otherwise $\widehat F_{ab}=\widetilde F_{ab}/\sigma_{ab}$
has a dimensionless spectrum.  We compare its adjacent-gap ratios, connected
form factor, and number variance with finite-dimensional reference ensembles.
For example, if $Z_F(\tau)=\Tr e^{-i\tau\widehat F_{ab}}$, then
\[
 K_{F,c}(\tau)=D^{-1}\left(
 \mathbb E|Z_F(\tau)|^2-|\mathbb EZ_F(\tau)|^2\right).
\]
The expectation is over the declared ensemble, not over eigenvalues of one
matrix.  Adjacent gaps within a common spectrum are correlated observations;
uncertainty intervals are therefore formed by resampling complete matrices or
complete registered panels.

Local repulsion and a finite-$D$ form-factor ramp test correlations within one
curvature component.  They do not test the joint distribution of different
tangent directions.  Nor do they determine whether deviations at longer
scales follow from a nontrivial covariance, a low-rank response algebra, or a
connected higher moment.  Those questions require matrix-element
contractions.

\subsection{Channel whitening and the four-point tensor}

Let $a=1,\ldots,M$ label a physical tangent family fixed before examining the
outcome.  At a common base projector, set
$Y_a=X_a-\mathbb E X_a$ and define its channel Gram matrix
\[
 G_{ab}=D^{-1}\mathbb E\Tr(Y_a^\dagger Y_b).
\]
We whiten on the registered support of $G$, using its Moore--Penrose inverse
when an exact null channel has been identified independently of the fourth
moment.  The resulting responses are denoted $\widehat X_a$.  Channel
whitening removes an overall choice of tangent coordinates, but it does not
assume that covariance factorizes between tangent, complement, and fiber
indices.

For the Moore--Read and lattice-supercharge panels, the mean is not fitted on
the inference sample.  It is evaluated exactly for the finite registered
tangent population: the uniform torus average over all guiding-center anchors
for Moore--Read, and the enumerated ordered local-coordinate marginals for the
lattice supercharge.  Panels 0--11 determine only the $2\times2$ channel
whitener; the frozen map is applied to panels 12--23, which form the primary
inference set.  The reversed split is recorded as a descriptive replication
and is not pooled with the primary estimate.

The gauge-invariant four-channel tensor is
\begin{equation}
 T_{abcd}=\frac{1}{D}\,
 \mathbb E\Tr\!\left(
 \widehat X_a^\dagger\widehat X_b
 \widehat X_c^\dagger\widehat X_d
 \right).
 \label{eq:fourpoint}
\end{equation}
A Gaussian comparison must reproduce the complete second moments of the
vectorized response.  If $z$ stacks all complement, fiber, and channel indices,
these moments are the covariance $C=\mathbb E(zz^\dagger)$ and the
pseudocovariance $\Pi=\mathbb E(zz^T)$.  The latter need not vanish for a
complex response ensemble.  We write their channel-indexed blocks as
$C_{ab}$ and $\Pi_{ab}$, with all ambient indices retained inside each block.
Isserlis' theorem therefore contains three Wick
pairings.  Suppressing the contracted ambient indices, their tensor structure
is
\[
 W_{abcd}=[C_{ab}C_{cd}]
 +[\Pi_{ac}^{*}\Pi_{bd}]
 +[C_{ad}C_{cb}].
\]
Keeping only the first and third terms imposes a proper-complex null that the
physical tangent ensemble need not satisfy.  Likewise, replacing each bracket
by a product of scalar channel covariances imposes an additional separability
assumption.  Neither restriction is part of the complete Gaussian comparison.

We define the connected fourth-order response by
\begin{equation}
 \mathcal K_{abcd}=T_{abcd}-W_{abcd},
 \qquad
 R_4^{\mathrm{full}}=\frac{\lVert\mathcal K\rVert_F}{\lVert W\rVert_F}.
 \label{eq:wickresidual}
\end{equation}
The ratio is reported only when $\lVert W\rVert_F$ passes the registered
nonzero-scale gate.  A scalar projection of $\mathcal K$ can be more powerful
than its norm, but its direction and sign must be fixed without reference to
the opened data.  A statistically resolved nonzero population contraction of
$\mathcal K$ rejects the matched Gaussian null for the stated ensemble and
tangent class; a nonzero finite-sample point estimate alone does not.  Such a
rejection does not identify a unique microscopic mechanism, and it does not
show that the residual survives as $D\rightarrow\infty$.

The registered scalar direction is tested in five primary cases.  We use the
one-sided Student-$t$ reference with $11$ degrees of freedom for the twelve
panel pseudovalues and control the family error at $0.05$ by a Bonferroni
threshold of $0.01$ per case.  This finite-sample gate, rather than the normal
approximation stored with the earlier raw-moment calculation, determines
whether a cell enters the claim table.

\subsection{Inference units and claim levels}

Every uncertainty statement names its declared statistical unit.  The
appropriate unit may be an independent disorder realization, an independently
drawn Hamiltonian, or a fixed-Hamiltonian tangent panel whose construction is
treated as exchangeable under a frozen protocol conditional on the common
base Hamiltonian.  These cases answer different
questions.  Eigenvalues and matrix elements drawn from the same response
matrix are not independent samples, and tangent directions evaluated at one
fixed projector do not estimate Hamiltonian-to-Hamiltonian fluctuations.
Pooling such quantities would understate uncertainty without adding an
independent physical realization.

For fixed-base panels, products entering the Wick tensor are estimated with
distinct panels, giving a U-statistic rather than a same-panel plug-in product.
Uncertainty is obtained from delete-one-panel pseudovalues, as detailed in the
Supplemental Material.  For disorder or independent-Hamiltonian ensembles,
the same construction is applied with the full realization as the deleted
unit.  Comparisons between models retain their native units unless a matched
sampling protocol has been carried out.

We consequently distinguish three levels of evidence.  Curvature
level repulsion is a finite-size statement about local eigenvalue
correlations.  Agreement of $T_{abcd}$ with the complete $W_{abcd}$ tests
fourth-order closure for a specified tangent ensemble.  Geometric ETH in an
asymptotic sense would require a growing-rank sequence, a fixed
operator class, and controlled decay of all covariance-whitened connected
cumulants.  The calculations in this paper are reported at the first two
levels only where their respective gates have actually been evaluated.
